\chapter{Document Elements}
\label{ch:elements}

This chapter describes the individual elements of which a thesis is
composed, from the typeface to the sectioning commands.

\section{Typefaces}
\label{sec:typefaces}

The document class requires the use of the ``Libertine'' typeface
family. Your \TeX\ installation should include this set of packages.
Please do not substitute other typefaces. The ``\verb|lmodern|'' and
``\verb|ltimes|'' packages should not be used, as they will override the
built-in typeface families.

\section{Title Information}
\label{sec:title}

The title of your work should use capital letters appropriately ---
\url{https://capitalizemytitle.com/} has useful rules for
capitalization. Set the title and subtitle with the standard
\verb|\title| command and the \verb|\subtitle| command:
\begin{verbatim}
  \title{The Name of the Title Is Hope}
  \subtitle{A Short Subtitle}
\end{verbatim}
Do not insert line breaks in your title.

\section{Metadata}
\label{sec:metadata}

The cover page is kept apart from the thesis class. It is provided by the
\texttt{cover} package, loaded and fed from the preamble:
\begin{verbatim}
  \usepackage{cover}
  \input{frontmatter/cover}
\end{verbatim}
The metadata file uses familiar ACM-style commands:
\begin{verbatim}
  \author{Ben Trovato}
  \affiliation{Department of Computer Science\\
               University of Somewhere}
  \email{trovato@corporation.com}
  \orcid{1234-5678-9012}
  \authornote{Corresponding author.}

  \degree{Doctor of Philosophy}
  \supervisor{Prof.\ Ada Lovelace}
  \cosupervisor{Dr.\ Alan Turing}
  \submissiondate{October 2026}
\end{verbatim}
Each value can be printed in the document with the corresponding
accessor, for example \verb|\theauthor|, \verb|\theaffiliation| and
\verb|\thesubmissiondate|.

The cover itself is produced by \verb|\maketitle|. To customise its
layout, redefine \verb|\maketitle| (or edit \texttt{cover.sty}) using
the accessors, for example:
\begin{verbatim}
  \renewcommand{\maketitle}{%
    \begin{titlepage}\centering
      {\Huge\bfseries\thetitle\par}\vspace{2cm}
      {\LARGE\theauthor\par}
    \end{titlepage}}
\end{verbatim}

\section{Sectioning Commands}
\label{sec:sectioning}

\sloppy
Your work should use standard \LaTeX\ sectioning commands:
\verb|\chapter|, \verb|\section|, \verb|\subsection|,
\verb|\subsubsection|, \verb|\paragraph|, and \verb|\subparagraph|. The
sectioning levels up to \verb|\subsubsection| should be numbered; do not
remove the numbering from the commands.
\fussy

Simulating a sectioning command by setting the first word or words of a
paragraph in boldface or italicized text is {\bfseries not allowed.}

Below are examples of sectioning commands.

\subsection{Subsection}
\label{sec:subsection}

This is a subsection.

\subsubsection{Subsubsection}
\label{sec:subsubsection}

This is a subsubsection.

\paragraph{Paragraph}

This is a paragraph.

\subparagraph{Subparagraph}

This is a subparagraph.

\section{Multi-language Documents}
\label{sec:languages}

Documents may be written in languages other than English, or include
titles, subtitles and abstracts in different languages. As a rule, a
document written in a language other than English should also include an
English title and abstract. The metadata commands accept arbitrary text,
so translated titles and abstracts can be set alongside the main ones.
